% !TEX program = lualatex
\documentclass[a4paper,11pt]{scrartcl}

\usepackage[
  lang=zh,
  mode=journal,
  theme=ocean,
  fontset=lm,
  extra-math=true,
  toc=true,
  toc-layout=page
]{synaptic}
\usepackage{booktabs,tabularx}
\newcolumntype{Y}{>{\raggedright\arraybackslash}X}

\SynapticNewTheorem[remark]{observation}{观察}
\SynapticHeader{synaptic 用户手册}
\SynapticSubHeader{面向 LuaLaTeX 的稳定学术排版}
\SynapticPreprint{v2.2.0}
\SynapticTitle{synaptic 宏包用户手册}
\SynapticAuthor{synaptic 项目}
\SynapticAffiliation{开源 LaTeX 社区}
\SynapticEmail{synaptic@example.com}
\SynapticSubject{synaptic 2.2.0 用户手册}
\SynapticKeywords{synaptic, LaTeX3, LuaLaTeX, 学术排版}

\begin{abstract}
  \textbf{synaptic} 是一套模块化 LaTeX3 学术文档设计系统，适用于期刊
  论文、书籍、讲义与笔记。本手册介绍它的环境要求、安全配置模型、元数据
  API、定理系统和各模式功能。
\end{abstract}

\begin{document}
\SynapticMakeTitle

\section{快速开始}

宏包要求 LuaLaTeX 和 KOMA-Script 文档类。仓库已包含生成好的宏包文件，
因此可直接使用本地副本：

\begin{verbatim}
export TEXINPUTS=/path/to/synaptic/tex/latex/synaptic//:$TEXINPUTS
lualatex document.tex
\end{verbatim}

最小论文示例：

\begin{verbatim}
\documentclass{scrartcl}
\usepackage[mode=journal]{synaptic}
\SynapticTitle{一篇短文}
\SynapticAuthor{作者}
\begin{abstract}摘要。\end{abstract}
\begin{document}
\SynapticMakeTitle
\section{引言}
正文。
\end{document}
\end{verbatim}

书籍模式应使用 \texttt{scrbook}。纸张尺寸由文档类决定，synaptic 不会强制 A4。

\section{配置}

\noindent\begin{tabularx}{\textwidth}{@{}l l Y@{}}
  \toprule
  \textbf{键} & \textbf{默认值} & \textbf{可选值与作用} \\
  \midrule
  \texttt{mode} & journal & journal、book、lecture 或 notes \\
  \texttt{theme} & ocean & ocean、graphite、forest、midnight 或 paper \\
  \texttt{lang} & en & en 或 zh，决定标签语言与中文支持 \\
  \texttt{fontset} & auto & auto、xcharter、libertinus、stix2 或 lm \\
  \texttt{toc} & false & 标题流程是否输出目录 \\
  \texttt{toc-layout} & top & top、inline 或 page \\
  \texttt{title-layout} & auto & inline、compact、page、cover 或模式默认值 \\
  \texttt{numbering} & auto & section、chapter、continuous、none 或模式默认值 \\
  \texttt{box-numbering} & shared & shared、separate 或 none \\
  \texttt{color-mode} & screen & screen、print 或 mono \\
  \texttt{density} & balanced & compact、balanced 或 airy \\
  \texttt{measure} & auto & narrow、standard、wide 或模式默认值 \\
  \texttt{binding-offset} & 0pt & 书籍装订补偿尺寸 \\
  \texttt{extra-math} & false & 便捷 Unicode 数学字母命令 \\
  \texttt{fine-breaking} & true & 适度放宽断行容忍度 \\
  \bottomrule
\end{tabularx}

仅 \texttt{theme}、\texttt{toc} 和 \texttt{toc-layout} 可在加载后安全修改：

\begin{verbatim}
\SynapticSetup{theme=paper,toc=false,toc-layout=page}
\end{verbatim}

模式、语言、字体与其他加载期排版选项会被冻结。运行时尝试修改时，宏包会给出
警告，而不是留下只完成一半的配置。

\subsection{模式}

\noindent\begin{tabularx}{\textwidth}{@{}l l Y@{}}
  \toprule
  \textbf{模式} & \textbf{建议文档类} & \textbf{用途} \\
  \midrule
  journal & scrartcl & 内联论文标题、元数据、运行页眉与声明 \\
  book & scrbook & 书籍扉页、篇章、镜像导航与前后文 \\
  lecture & scrartcl/scrreprt & 课程信息条与语义教学系统 \\
  notes & scrartcl & 紧凑信息条与轻量知识卡 \\
  \bottomrule
\end{tabularx}

\subsection{字体}

\texttt{auto} 依次选择第一套完整安装的组合：XCharter、Libertinus、STIX Two、
Latin Modern。显式选择字体组时，如果缺少必要的配套字体，宏包会清晰报错。
中文优先使用 Noto CJK SC，并以 TeX Live 自带的 Fandol 作为后备。

\section{标题与 PDF 元数据}

\noindent\begin{tabularx}{\textwidth}{@{}l Y@{}}
  \toprule
  \textbf{命令} & \textbf{含义} \\
  \midrule
  \SynapticCmd{SynapticTitle} & 必填标题与 PDF 标题 \\
  \SynapticCmd{SynapticSubtitle} & 可选副标题 \\
  \SynapticCmd{SynapticShortTitle} & 运行页眉短标题 \\
  \SynapticCmd{SynapticDate} & 显式日期 \\
  \SynapticCmd{SynapticAuthor} & 作者及可选机构标记，可重复 \\
  \SynapticCmd{SynapticAffiliation} & 机构，可重复 \\
  \SynapticCmd{SynapticEmail} & 可点击联系方式，可重复 \\
  \SynapticCmd{SynapticSubject} & PDF 主题 \\
  \SynapticCmd{SynapticKeywords} & 页面关键词与 PDF 关键词 \\
  \SynapticCmd{SynapticHeader} & 投稿抬头 \\
  \SynapticCmd{SynapticSubHeader} & 次级抬头 \\
  \SynapticCmd{SynapticPreprint} & 预印本编号 \\
  \SynapticCmd{SynapticArxiv} & arXiv 链接 \\
  \SynapticCmd{SynapticDedicated} & 献词 \\
  \SynapticCmd{SynapticMakeTitle} & 验证并输出标题信息 \\
  \bottomrule
\end{tabularx}

四种模式共享元数据，但采用不同标题渲染器。journal 支持出版信息、通讯信息与作者注；
book 支持半扉页、出版社、版次、ISBN 与版权；lecture 支持课程、课程代码、讲次、教师与
学期。论文的基金、数据、伦理或利益冲突声明可用
\verb|\SynapticStatement{标题}{正文}|，其星号形式不写入目录。未设置 Synaptic 字段时，
\verb|\title|、\verb|\author| 和 \verb|\date| 可被导入。
PDF 字符串只进行一次 Unicode 安全转换。

\section{定理与数学}

内置 \texttt{theorem}、\texttt{lemma}、\texttt{proposition} 与 \texttt{definition}
共享一个计数器。书籍默认按章编号，短文档默认按节编号。定义正文使用正体；中文命题类
正文保持正体，避免伪斜体造成的阅读干扰。

\begin{theorem}[高斯积分]
  \[
    \int_0^\infty e^{-x^2}\,\mathrm{d}x=\frac{\sqrt{\pi}}{2}.
  \]
\end{theorem}

\begin{proof}
  对积分求平方，转换为极坐标，再计算径向积分。
\end{proof}

自定义定理可选择 \texttt{plain}、\texttt{definition} 或 \texttt{remark} 样式：

\begin{verbatim}
\SynapticNewTheorem[remark]{observation}{观察}
\end{verbatim}

\begin{observation}
  开启 \texttt{extra-math=true} 后可使用 \(\symbb{R}\)、\(\symcal{F}\) 与
  \(\symfrak{g}\) 等命令。
\end{observation}

\section{主题与讲义盒子}

使用 \verb|\SynapticUseTheme{名称}| 切换主题。下列命令使用 xcolor 颜色表达式
定义并立即启用自定义主题：

\begin{verbatim}
\SynapticDefineTheme{brand}{blue!70!black}{gray!80!black}[orange]
\end{verbatim}

讲义模式增加 \texttt{synexample}、\texttt{synremark}、\texttt{synwarning}、
\texttt{synexercise} 和 \texttt{synbox}。例题与练习默认共享定理计数器，注记与警告
默认不编号；第一个可选参数是标题，第二个是标签。课程单元还可使用
\texttt{synlearninggoals} 与 \texttt{synlecturesummary}。

\section{书籍流程与致谢}

书籍中可使用 \SynapticCmd{SynapticFrontMatter}、\SynapticCmd{SynapticMainMatter} 和
\SynapticCmd{SynapticBackMatter}，以及目录、图表目录、附录、章副标题与题词辅助命令。
notes 模式提供 \texttt{synkeyidea}、\texttt{synquestion}、\texttt{syntodo} 和
\texttt{synsummary}，并兼容文档类的双栏布局。\SynapticCmd{SynapticAcknowledgments} 在书籍模式中生成章，
其他模式中生成节；带星号的形式不加入目录。

\section{开发命令}

\begin{verbatim}
l3build unpack   # 从 synaptic.dtx 重新生成宏包文件
l3build check    # 运行 LuaLaTeX 回归测试
l3build doc      # 生成源码文档与手册
l3build ctan     # 创建发布压缩包
\end{verbatim}

\SynapticAcknowledgments*
项目感谢 LaTeX3、KOMA-Script、fontspec 与 TeX 社区。

\end{document}
